\chapter*{Prefacio}
\addcontentsline{toc}{chapter}{Prefacio}

\begin{flushright}
\itshape
``The underlying physical laws necessary for the mathematical theory of a\\
large part of physics and the whole of chemistry are thus completely known,\\
and the difficulty is only that the exact application of these laws leads to\\
equations much too complicated to be soluble.''\\[4pt]
\upshape --- P.\,A.\,M.\ Dirac (1929)
\end{flushright}
\bigskip

Dirac tenía razón: las leyes se conocen, pero las ecuaciones son intratables.
Noventa y cinco años después, la situación es la misma y distinta. Las
ecuaciones siguen siendo intratables —ningún computador resolverá la ecuación
de Schrödinger para $10^{23}$ partículas—, pero ha aparecido una clase de
herramientas que no intentan resolverlas: las \emph{aprenden}. Una red neuronal
no integra la ecuación de Navier-Stokes; aproxima el mapa entre condiciones
iniciales y soluciones a partir de ejemplos. Un modelo de difusión no resuelve
la ecuación de Fokker-Planck; la invierte estocásticamente para generar muestras
de distribuciones complejas. La pregunta que organiza este libro es: ¿qué
tienen estas herramientas en común con la física que el lector ya conoce?

La respuesta es: todo. Las redes neuronales, el descenso de gradiente
estocástico, los modelos de difusión y el transporte óptimo no son
herramientas computacionales importadas a la física desde la informática. Son
física. La ecuación de Fokker-Planck que describe la distribución estacionaria
del SGD es la misma que describe la distribución de neutrones en un medio
difusivo. El \textit{score} $\nabla_{\bm{x}}\log p_t$ que guía la generación
de imágenes es la misma cantidad que aparece en la condición de balance
detallado de la mecánica estadística de equilibrio. El puente de Schrödinger
que unifica los modelos de difusión es el mismo problema variacional que
Schrödinger formuló en 1931 pensando en partículas brownianas, no en
imágenes. Estas no son analogías: son identidades matemáticas. El
Apéndice~\ref{ap:info_mecanica} las presenta todas juntas en una sola tabla.

\subsection*{El lector al que está dirigido este libro}

Este libro está escrito para un estudiante de posgrado en física ---o de
último año de licenciatura con motivación suficiente--- que ha tomado los
cursos estándar de mecánica clásica, electromagnetismo, mecánica cuántica,
mecánica estadística y métodos matemáticos. Se asume familiaridad con
ecuaciones diferenciales ordinarias y parciales, álgebra lineal y cálculo
multivariable. Se asume que el lector sabe programar en algún lenguaje
científico, aunque nunca ha implementado una red neuronal.

No se asume conocimiento de física nuclear ni de física de altas energías.
Los conceptos de estos dominios aparecen únicamente en las secciones
opcionales descritas más adelante, y siempre con definición explícita.
La complejidad matemática escala gradualmente; cada sección indica su
nivel de exigencia en dos ejes independientes: dificultad matemática y
dificultad física, representados mediante estrellas al inicio de cada
sección.

\subsection*{Estructura del libro y cómo leerlo}

El libro se organiza en doce capítulos que forman tres arcos encadenados.
Un ejemplo unificado ---el ajuste de la señal de un oscilador armónico
simple--- recorre los primeros cuatro capítulos como laboratorio práctico,
permitiendo al lector comparar cada herramienta nueva contra las anteriores
en un sistema que conoce de memoria.

Los \textbf{Capítulos~1--4} construyen el problema de aproximación de
funciones y las herramientas para resolverlo. El Capítulo~1 parte de la
pregunta que abre el libro ---cómo aproximar una función desconocida a
partir de datos finitos--- y desarrolla las familias paramétricas, las
ecuaciones normales (como el equilibrio de un resorte en $\mathbb{R}^p$),
y las funciones de base radial como puente entre la interpolación clásica
y las redes neuronales. El Capítulo~2 formaliza el compromiso
sesgo-varianza, introduce la regularización de Tikhonov, y revela mediante
la SVD que la regularización es un filtro espectral ---idéntico al filtro de
Wiener de la teoría de señales. El Capítulo~3 estudia la geometría del
paisaje de pérdida: gradiente, hessiano, índice de Morse, convexidad, y la
observación (apoyada en la teoría de matrices aleatorias) de que en alta
dimensión los mínimos locales son exponencialmente raros. El Capítulo~4
introduce el perceptrón multicapa como familia de funciones con bases
adaptativas, conecta las funciones de activación con la física
(la sigmoide es la distribución de Fermi-Dirac, la tanh es la función de
Langevin paramagnética), y demuestra el teorema de aproximación universal.

Los \textbf{Capítulos~5--9} introducen el ruido, la incertidumbre y la
estructura. El Capítulo~5 establece la conexión central del libro: el SGD
con ruido de mini-lote es una ecuación de Langevin, y su distribución
estacionaria es la distribución de Gibbs de la mecánica estadística. El
Capítulo~6 desarrolla la backpropagation como caso discreto del método
adjunto de Pontryagin. El Capítulo~7 provee el fundamento probabilístico:
la pérdida MSE es el MLE gaussiano, la regularización $\ell_2$ es un prior
gaussiano, y la ELBO es la energía libre variacional de Helmholtz. Los
Capítulos~8 y~9 completan el arsenal: regularización práctica, normalización,
el fenómeno \textit{double descent}, y las arquitecturas especializadas
(CNNs como equivarianza de traslación, GNNs, transformers, redes
equivariantes de grupo con coeficientes de Clebsch-Gordan, Neural ODEs).

Los \textbf{Capítulos~10--12} culminan el arco. El Capítulo~10 construye
el cálculo de Itô con rigor: la variación cuadrática del movimiento
browniano, la integral de Itô como límite $L^2$, y la fórmula de Itô como
regla de la cadena estocástica. El Capítulo~11 deriva la ecuación de
Fokker-Planck por dualidad $L^2$, demuestra el teorema de Anderson sobre
la reversión temporal de SDEs, desarrolla el score matching como pérdida
computable, y conecta los modelos de difusión con el puente de Schrödinger
y el transporte óptimo. El Capítulo~12 ensambla todo en algoritmos
concretos: DDPM, score SDE, DDIM, flow matching, OT flow matching y
classifier-free guidance, con las arquitecturas U-Net y DiT en el espacio
latente.

\subsection*{Los apéndices}

El libro incluye cinco apéndices diseñados para hacer el texto
autocontenido sin sobrecargar los capítulos. El
\textbf{Apéndice~\ref{ap:probabilidad}} cubre los conceptos de
probabilidad que van más allá del currículo estándar de física: filtración,
esperanza condicional como proyección $L^2$, martingalas y el lema de
Borel-Cantelli. El \textbf{Apéndice~\ref{ap:info_mecanica}} es el argumento
central del libro en forma comprimida: la identidad entre entropía de
Boltzmann y entropía de Shannon, la distribución de Gibbs como máxima
entropía, el H-teorema como decrecimiento de la KL, y la tabla completa
de correspondencias física estadística--aprendizaje automático. El
\textbf{Apéndice~\ref{ap:algebra_lineal}} cubre la SVD, la pseudoinversa,
las normas matriciales, el hessiano y el espacio $L^2$ como espacio de
Hilbert. El \textbf{Apéndice~\ref{ap:algoritmos}} consolida los algoritmos
de implementación: Euler-Maruyama, Milstein, el bucle de entrenamiento por
denoising score matching, y una guía de referencia de las funciones de
NumPy, PyTorch y JAX. El \textbf{Apéndice~\ref{ap:glosario}} es un
diccionario de correspondencias entre el vocabulario del aprendizaje
automático y su análogo físico.

\subsection*{La Estrella Polar: un hilo narrativo opcional}

Al final de cada capítulo, después de los ejercicios, el lector encontrará
una sección independiente llamada \textbf{Estrella Polar}. Estas secciones
son completamente opcionales: pueden omitirse sin perder ningún concepto del
texto principal ni interrumpir la continuidad entre capítulos. Cada Estrella
Polar es autocontenida y no se mezcla con la narrativa del capítulo que la
precede.

Las Estrellas Polares aplican los resultados de cada capítulo a sistemas
físicos concretos. Los Capítulos~1 a~8 las usan para la predicción del
flujo neutrónico en el reactor de investigación IAN-R1: un sistema de fisión
con seis grupos de neutrones precursores donde los conceptos del libro ---la
ecuación de Langevin, el balance detallado, el filtrado bayesiano, la
ecuación de Fokker-Planck con coeficientes variables--- tienen
interpretación física inmediata y parámetros verificables
experimentalmente. A partir del Capítulo~9, un segundo sistema aparece
junto al reactor: la simulación de cascadas electromagnéticas en
calorímetros de física de altas energías. En el Capítulo~12, ambos
sistemas se integran en una narrativa unificada que demuestra que las
mismas herramientas matemáticas se transfieren entre dominios físicos sin
modificación.

Para el lector que las siga, las Estrellas Polares construyen una narrativa
paralela y autocontenida. Para el lector que las omita, el libro es un
texto autónomo de procesos estocásticos y aprendizaje automático para
físicos, sin dependencia alguna de la física nuclear ni de la física de
altas energías.

\subsection*{Los anglicismos técnicos}

Los términos \textit{score matching}, \textit{flow matching},
\textit{dropout}, \textit{batch normalization} y similares se conservan
en inglés. Un físico que acceda a la literatura científica actual necesita
reconocerlos. Cuando existe una traducción consolidada se usa; cuando no,
se emplea el término original con definición explícita la primera vez que
aparece.

\subsection*{Convenciones notacionales}

\begin{itemize}
  \item Escalares: letras minúsculas en itálica ($x,\, t,\, \lambda$).
  \item Vectores columna: letras minúsculas en negrita
        ($\bm{x},\, \bm{n},\, \bm{\phi}$).
  \item Matrices y tensores: letras mayúsculas en negrita
        ($\bm{D},\, \bm{\Sigma},\, \bm{J}$).
  \item Conjuntos y espacios de funciones: letras caligráficas
        ($\mathcal{L},\, \mathcal{A},\, \mathcal{N}$).
  \item Esperanza: $\langle\,\cdot\,\rangle$ en contexto físico,
        $\mathbb{E}[\,\cdot\,]$ en contexto probabilístico. Ambas
        notaciones son intercambiables.
  \item Procesos estocásticos: $\bm{X}_t$ para el proceso como objeto
        matemático, $\bm{x}$ para una realización particular.
  \item Función de \textit{score}:
        $\bm{s}(\bm{x},t) = \nabla_{\bm{x}}\log p(\bm{x},t)$.
  \item Reactividad nuclear (solo en secciones Estrella Polar):
        $\varrho$, para distinguirla de la densidad de masa $\rho$.
\end{itemize}

\subsection*{Sobre el título}

\textit{La máquina estocástica} no es una metáfora. El universo procesa
información: un sistema de fisión computa cuántos neutrones sobreviven cada
generación, un proceso de difusión computa qué configuraciones son
termodinámicamente probables, una red neuronal computa qué distribución de
parámetros minimiza una energía libre. Ninguna de estas computaciones es
determinista. Todas están gobernadas por las mismas ecuaciones. Este libro
estudia esas ecuaciones.

\vspace{1.2cm}
\begin{flushright}
\textit{J.\,A. Chala-Casanova}\\[2pt]
\the\year
\end{flushright}